\documentclass[11pt]{article}
\usepackage[a4paper,margin=1in]{geometry}
\usepackage{amsmath,amssymb}
\usepackage[T1]{fontenc}
\usepackage{lmodern}
\usepackage{microtype}
\usepackage{xcolor}

\setlength{\parindent}{0pt}
\setlength{\parskip}{0.5em}

\newcommand{\promptbox}[1]{%
  \begingroup\setlength{\fboxsep}{10pt}\noindent
  \colorbox{blue!8}{\parbox{\dimexpr\linewidth-2\fboxsep\relax}{\textbf{Prompt. }#1}}
  \endgroup
}
\newcommand{\resultbox}[1]{%
  \begingroup\setlength{\fboxsep}{10pt}\noindent
  \colorbox{purple!8}{\parbox{\dimexpr\linewidth-2\fboxsep\relax}{\textbf{Result. }#1}}
  \endgroup
}
\newcommand{\examplebox}[1]{%
  \begingroup\setlength{\fboxsep}{10pt}\noindent
  \colorbox{green!8}{\parbox{\dimexpr\linewidth-2\fboxsep\relax}{\textbf{Example. }#1}}
  \endgroup
}

\title{\textbf{Book 1 --- Lecture 1 --- Interlude 1 --- Exercise 6}\\[2mm]
\large Why Orthogonal Vectors Have Zero Dot Product}
\author{}
\date{}

\begin{document}
\maketitle

\promptbox{Can you explain why the dot product of two vectors that are
orthogonal is \(0\)?}

\section*{Geometric explanation}
The dot product can be written as
\[
\mathbf{A}\cdot\mathbf{B} = \lVert\mathbf{A}\rVert\lVert\mathbf{B}\rVert\cos\theta,
\]
where \(\theta\)\footnote{\(\theta\) is the angle between the vectors.} is the
angle between the vectors.

If vectors are orthogonal, then \(\theta=90^\circ\) (or \(\pi/2\) radians), and
\(\cos(90^\circ)=0\). Therefore,
\[
\mathbf{A}\cdot\mathbf{B}
= \lVert\mathbf{A}\rVert\lVert\mathbf{B}\rVert\cdot 0
= 0.
\]

\section*{Component explanation}
In coordinates, dot product is
\[
\mathbf{A}\cdot\mathbf{B}=\sum_i A_i B_i.
\]
Orthogonality means there is no component of one vector along the direction of
the other\footnote{Equivalently, the projection of one onto the other has zero
length.}, so the signed component products cancel to zero.

\resultbox{Orthogonal vectors have zero dot product because the geometric factor
\(\cos\theta\) becomes zero at \(90^\circ\).}

\examplebox{In 2D, \((1,0)\) and \((0,5)\) are perpendicular. Their dot product
is \(1\cdot0+0\cdot5=0\).}

\end{document}
